\section{Phương trình mặt cầu}
\subsection{Đề 1}
\Opensolutionfile{ans}[ans/ans-2-B17-De1]
\caulc
\begin{ex}%[2H5N3-1]
	Trong không gian với hệ trục tọa độ $Oxyz$, phương trình nào sau đây là phương trình mặt cầu
	\choice
	{\True $x^2+y^2+z^2-2x=0$}
	{$x^2+y^2-z^2+2x-y+1=0$}
	{$2x^2+2y^2=\left( x+y \right)^{2}-z^2+2x-1 $}
	{$\left(x+y \right)^{2}=2xy-z^2-1 $}
	\loigiai{Phương trình mặt cầu $\left(S \right)$ có hai dạng là:
\begin{description}
\item[{\rm (1)}] $\left(x-a \right)^2+\left(y-b \right)^2+\left(z-c \right)^2=R^2$
\item[{\rm (2)}] $x^2+y^2+z^2-2ax-2by-2cz+d=0 \text{ với } a^2+b^2+c^2-d>0$
\end{description}	
$\ast$ $x^2+y^2+z^2-2x=0$ là phương trình mặt cầu.\\
$\ast$ $x^2+y^2-z^2+2x-y+1=0$ không là phương trình mặt cầu.\\
$\ast$ $2x^2+2y^2=\left( x+y \right)^{2}-z^2+2x-1 \Leftrightarrow x^2+y^2+z^2-2xy-2x+1=0$ không là phương trình mặt cầu.\\
$\ast$ $\left(x+y \right)^{2}=2xy-z^2-1 \Leftrightarrow x^2+y^2+z^2+1=0$ không là phương trình mặt cầu vì $a^2+b^2+c^2-d=-1<0$.  
}
\end{ex}	

\begin{ex}%[2H5N3-2]
Trong không gian với hệ trục tọa độ $Oxyz$, mặt cầu $\left(S \right) \left( x-1\right)^2+y^2+\left( z+3\right)^2=16$ có tâm là
	\choice
	{$I\left(1;0;3\right)$}
	{$I\left(-1;0;-3\right)$}
	{\True $I\left(1;0;-3\right)$}
	{$I\left(1;2;-3\right)$}
	\loigiai{
		Phương trình mặt cầu $\left( S \right)$ có dạng là:
$\left(x-a \right)^2+\left(y-b \right)^2+\left(z-c \right)^2=R^2$ có tâm $I\left(a;b;c\right)$ bán kính là $R$\\
Khi đó mặt cầu $\left(S\right)$ có tâm $I\left(1;0;-3\right)$.}
\end{ex}

\begin{ex}%[2H5N3-2]
Trong không gian với hệ trục tọa độ $Oxyz$, mặt cầu $x^2+y^2+z^2-4x+1=0$ có tâm và bán kính là
	\choice
	{$I\left(2;0;0\right)$, $R=3$}
	{$I\left(-2;0;0\right)$, $R=\sqrt{3}$}
	{$I\left(0;2;0\right)$, $R=\sqrt{3}$}
	{\True $I\left(2;0;0\right)$, $R=\sqrt{3}$}
	\loigiai{
		Phương trình mặt cầu $\left( S \right)$ có dạng là:
$x^2+y^2+z^2-2ax-2by-2cz+d=0$ có tâm $I\left(a;b;c\right)$, có bán kính là $R =\sqrt{a^2+b^2+c^2-d}$\\
Khi đó, mặt cầu $x^2+y^2+z^2-4x+1=0$ có tâm là $I\left(2;0;0\right)$, có bán kính là $R=\sqrt{2^2+0^2+0^2-1}=\sqrt{3}$.}
\end{ex}
\begin{ex}%[2H5N3-3] 
Trong không gian với hệ trục tọa độ $Oxyz$, phương trình mặt cầu có tâm $I\left(-3;1;2\right)$, bán kính $R=3$ là
	\choice
	{$\left(x+3 \right)^2+\left(y-1 \right)^2+\left(z+2 \right)^2=3$}
	{$\left(x-1 \right)^2+\left(y+3 \right)^2+\left(z-2 \right)^2=9$}
	{$\left(x-3 \right)^2+\left(y+1 \right)^2+\left(z+2 \right)^2=9$}
	{\True $\left(x+3 \right)^2+\left(y-1 \right)^2+\left(z-2 \right)^2=9$}
	\loigiai{
		Phương trình mặt cầu có tâm $I\left(-3;1;2\right)$, có bán kính $R=3$ là: $$\left(x+3 \right)^2+\left(y-1 \right)^2+\left(z-2 \right)^2=9$$ }
\end{ex}
\begin{ex}%[2H5H3-2]
Mặt cầu ${3x}^2+{3y}^2+{3z}^2-6x+12y+2=0$ có bán kính bằng
\choice
	{\True $R=\sqrt{\dfrac{13}{3}}$}
	{$R=\dfrac{\sqrt{7}}{3}$}
	{$R=\dfrac{2\sqrt{7}}{3}$}
	{$R=\dfrac{\sqrt{21}}{3}$}
	\loigiai{
		Biến đổi ${3x}^2+{3y}^2+{3z}^2-6x+12y+2=0 \Leftrightarrow x^2+y^2+z^2-2x+4y+\dfrac{2}{3}=0$ có tâm $I\left(1;-2;0\right)$ và bán kính $R= \sqrt{1^2+\left(-2 \right)^2+0^2-\dfrac{2}{3}}=\sqrt{\dfrac{13}{3}}$.}
\end{ex}
\begin{ex}%[2H5H3-1]
Mặt cầu $\left(S \right)$ : $x^2+y^2+z^2-2x+10y+3z+1=0$ đi qua điểm có tọa độ nào sau đây
\choice
	{$\left(3;-2;-4 \right) $}
	{\True $\left(4;-1;0 \right) $}
	{$\left(2;1;9 \right) $}
	{$\left(-1;3;-1 \right) $}
	\loigiai{
	Lần lượt thay các tọa độ các điểm vào phương trình mặt cầu. Tọa độ điểm nào thỏa mãn phương trình thì điểm đó thuộc mặt cầu.\\
	$\ast$ Với đáp án $\textbf{A}$, ta thay $x=3,y=-2,z=-4$ vào $\left(S\right)$.\\
	 Ta có $3^2+\left(-2\right)^2+\left(-4\right)^2-2.3+10.\left(-2\right)+3.\left(-4\right)+1=-8\neq0$. Khi đó $\left(3;-2;-4\right) \notin \left(S\right)$ nên ta loại đáp án $\textbf{A}$\\
	$\ast$ Tương tự với đáp án $\textbf{B}$. Ta có $4^2+\left(-1\right)^2+0^2-2.4+10.\left(-1\right)+3.0+1=0$. Khi đó $\left(4;-1;0\right) \in \left(S\right)$ nên ta chọn $\textbf{B}$\\
	$\ast$ Với đáp án $\textbf{C}$. Ta có $2^2+1^2+9^2-2.2+10.1+3.9+1=120$. Khi đó $\left(2;1;9\right) \notin \left(S\right)$ nên ta loại đáp án $\textbf{C}$\\
	$\ast$ Với đáp án $\textbf{D}$. Ta có $\left(-1\right)^2+3^2+\left(-1\right)^2-2.\left(-1\right)+10.3+3.\left(-1\right)+1=41$. Khi đó $\left(-1;3;-1\right) \notin \left(S\right)$ nên ta loại đáp án $\textbf{D}$}
\end{ex}
\begin{ex}%[2H5H3-3]
Mặt cầu tâm $I\left(-1;2;-3 \right)$ và đi qua điểm $A\left(2;0;0 \right)$ có phương trình là
\choice
	{\True $\left(x+1 \right)^2+\left(y-2 \right)^2+\left(z+3 \right)^2=22$}
	{$\left(x+1 \right)^2+\left(y-2 \right)^2+\left(z+3 \right)^2=11$}
	{$\left(x-1 \right)^2+\left(y+2 \right)^2+\left(z-3 \right)^2=22$}
	{$\left(x-1 \right)^2+\left(y-2 \right)^2+\left(z-3 \right)^2=22$}
	\loigiai{
	Ta có $\vec{IA}=\left( 3;-2;3\right)$, $R=IA=\sqrt{22}$\\
	Vậy phương trình mặt cầu tâm $I\left(-1;2;-3 \right)$, bán kính $R=\sqrt{22}$ là: $$\left(x+1 \right)^2+\left(y-2 \right)^2+\left(z+3 \right)^2=22$$
	}
\end{ex}
\begin{ex}%[2H5H3-2]
Mặt cầu $\left(S \right):\left(x+y \right)^2=2xy-z^2+1-4x$ có tâm là
\choice
	{$I\left(4;0;0\right)$}
	{\True $I\left(-2;0;0\right)$}
	{$I\left(-4;0;0\right)$}
	{$I\left(2;0;0\right)$}
	\loigiai{
	Ta biến đổi $\left(x+y\right)^2=2xy-z^2+1-4x \Leftrightarrow x^2+y^2+z^2+4x-1=0$\\
	Vậy mặt cầu có tâm $I\left(-2;0;0\right)$.}
\end{ex}
\begin{ex}%[2H5H3-3]
Trong không gian với hệ trục tọa độ $Oxyz$, cho hai điểm $A\left(1;0;-3 \right)$ và $B\left(3;2;1 \right)$. Phương trình mặt cầu đường kính $AB$ là
\choice
	{$x^2+y^2+z^2+4x-2y+2z=0$}
	{\True $x^2+y^2+z^2-4x-2y+2z=0$}
	{$x^2+y^2+z^2-2x-y+z-6=0$}
	{$x^2+y^2+z^2-4x-2y+2z+6=0$}
	\loigiai{
Ta có $\vec{AB}=\left(2;2;4\right) \Rightarrow AB=2\sqrt{6}$.
	Mặt cầu đường kính $AB$ có tâm $I$ là trung điểm $AB$ nên $I\left(2;1;-1\right)$, bán kính $R=\dfrac{AB}{2}=\sqrt{6}$.\\
	Khi đó phương trình mặt cầu: $\left(x-2\right)^2+\left(y-1\right)^2+\left(z+1\right)^2=6$.\\
	 hay $x^2+y^2+z^2-4x-2y+2z=0$.}
\end{ex}
\begin{ex}%[2H5V3-3]
Phương trình mặt cầu $\left(S \right)$ có tâm $O$, tiếp xúc với mặt phẳng $\left(\alpha\right):16x-15y-12z+75=0$ là
\choice
	{$x^2+y^2+z^2-3x=9$}
	{$x^2+y^2+z^2=3$}
	{$x^2+y^2+z^2=81$}
	{\True $x^2+y^2+z^2=9$}
	\loigiai{
	Do $\left( S\right)$ tiếp xúc với $\left(\alpha\right)$ nên $\mathrm{d}\left(O;\left( \alpha\right)\right) = R \Leftrightarrow R = \dfrac{75}{25}=3$.\\
	Mặt cầu tâm $O\left( 0;0;0\right)$ và bán kính $R=3$ có phương trình $\left( S\right): x^2+y^2+z^2=9$.}
\end{ex}
\begin{ex}%[2H5V3-3]
Cho mặt phẳng $\left(P\right)$ và mặt cầu $\left(S\right)$ có phương trình lần lượt là\\
 $\left(P\right): 2x+2y+z-m^2+4m-5=0$ và $\left(S\right): x^2+y^2+z^2-2x+2y-2z-6=0$. Giá trị của $m$ để $\left(P\right)$ tiếp xúc với $\left(S\right)$ là
\choice
	{\True $m=-1$ hoặc $m=5$}
	{$m=1$ hoặc $m=-5$}
	{$m=-1$}
	{$m=5$}
	\loigiai{
	Dễ thấy mặt cầu $\left(S \right)$ có tâm $I\left(1;-1;1\right)$ và bán kính $R=\sqrt{1^2+\left( -1\right)^2+1^2+6}=3$.\\
	Khi $\left(P\right)$ tiếp xúc với $\left(S\right)$ ta có:
\begin{eqnarray*}
& &\mathrm{d}\left(I;\left(P \right)\right)=R\\
&\Leftrightarrow & \dfrac{\vert 2\cdot1+2\cdot\left(-1\right)+1\cdot1-m^2+4m-5 \vert}{\sqrt{2^2+2^2+1^2}}=3\\
&\Leftrightarrow & \vert m^2-4m+4 \vert = 9\\
&\Leftrightarrow & \hoac{&m^2-4m+4=9\\&m^2-4m+4=-9}\\
&\Rightarrow & m^2-4m-5=0\\
&\Rightarrow &\hoac{&m=-1\\&m=5}.
\end{eqnarray*}}
\end{ex}
\begin{ex}%[2H5V3-3]
Cho bốn điểm $A\left(6;-2;3\right)$,$B\left(0;1;6\right)$,$C\left(2;0;-1 \right)$,$D\left(4;1;0\right)$. Khi đó mặt cầu ngoại tiếp tứ diện $ABCD$ có phương trình là
\choice
	{$x^2+y^2+z^2-2x+y-3z-3=0$}
	{\True $x^2+y^2+z^2-4x+2y-6z-3=0$}
	{$x^2+y^2+z^2+4x-2y+6z-3=0$}
	{$x^2+y^2+z^2+2x-y+3z-3=0$}
	\loigiai{
	Gọi phương trình mặt cầu $\left( S\right)$ có dạng: $x^2+y^2+z^2-2Ax-2By-2Cz+D=0$ với $A^2+B^2+C^2-D>0$, ta có\\
	$\heva{&A\left(6;-2;3\right)\in \left( S\right)\\&B\left(0;1;6\right)\in \left( S\right)\\&C\left(2;0;-1\right)\in \left( S\right)\\&D\left(4;1;0\right)\in \left( S\right)}$ $\Leftrightarrow \heva{&49-12A+4B-6C+D=0\quad(1)\\&37   -2B-12C+D=0\quad(2)\\&5-4A    +2C+D=0\quad(3)\\&17-8A-2B   +D=0\quad(4)}$\\
	Lấy $\left(1 \right)-\left( 2\right)$;$\left(2 \right)-\left( 3\right)$;$\left(3 \right)-\left( 4\right)$ ta được hệ\\
	$\heva{&-12A+6B+6C=-12\\&4A-2B-14C=-32\\&4A+2B+2C=12} \Leftrightarrow \heva{&A=2\\&B=-1\\&C=3} \Rightarrow D=3$.\\
	Vậy phương trình mặt cầu là: $x^2+y^2+z^2-4x+2y-6z-3=0$.
	
	\medskip
	Ở bài này máy tính Casio cũng có thể thực hiện giải hệ phương trình bậc nhất bốn ẩn để tìm các hệ số.
	}
\end{ex}
\Closesolutionfile{ans}
\indapan{6}{ans/ans-2-B17-De1}

\cauds
\Opensolutionfile{ans}[ans/ans-0-B15-DS]
\begin{ex}%[2H5H3-2]
Cho $\left(S\right): x^2+y^2+z^2-2x-4y+6z-67=0$. Các mệnh đề sau đúng hay sai ?
\choiceTF
	{Mặt cầu $\left( S\right)$ có tâm $I\left(1;2;3\right)$}
	{\True Bán kính mặt cầu $\left( S\right)$ là $R=9$}
	{Cho mặt phẳng $\left(P\right): 2x-2y+z-13=0$. Khi đó $\left( P\right)$ tiếp xúc với $\left(S\right)$}
	{\True Cho đường thẳng $\left(\Delta\right) :\heva{&x=1+t\\&y=2\\&z=-4+7t}$. Khi đó $\left(\Delta\right)$ và $\left(S\right)$ cắt nhau tại hai điểm}
	\loigiai{
		\begin{itemchoice}
			\itemch Sai. Mặt cầu $\left( S\right)$ có tâm $I\left(1;2;-3\right)$.
			\itemch Đúng. Ta có $R=\sqrt{1^2+2^2+\left(-3\right)^2+67}=\sqrt{81}=9$
			\itemch Sai. Vì $\mathrm{d}\left(I,\left(P\right)\right)=\dfrac{\vert 2\cdot1-2\cdot2+\left(-3\right)-13\vert}{\sqrt{2^2+2^2+\left(-1\right)^2}}=6< R=9$. Suy ra $\left(P\right)$ cắt $\left(S\right)$ theo giao tuyến là một đường tròn.
			\itemch Đúng. Tọa độ giao điểm là nghiệm của hệ phương trình\\
	 $\heva{&x=1+t\\&y=2\\&z=-4+7t\\&x^2+y^2+z^2-2x-4y+6z-67=0}\Rightarrow  \hoac{&t=0\\&t=1}$.\\
	 Với $t=0$, thay $t=0$ vào $\left(\Delta\right)$ ta được tọa độ $A\left(1;2;-4\right)$\\
	 Với $t=1$, thay $t=1$ vào $\left(\Delta\right)$ ta được tọa độ $B\left(2;2;3\right)$.\\
	 Vậy, $\left(\Delta\right)$ và $\left(S\right)$ cắt nhau tại hai điểm $A$ và $B$.
		\end{itemchoice}
	}
\end{ex}
\begin{ex}%[2H5H3-3]
	Trong không gian $Oxyz$ cho ba điểm $A\left(1;0;0\right)$ $B\left(0;2;0\right)$ $C\left(0;0;3\right)$ và phương trình mặt cầu $\left(S \right): \left(x-3\right)^2+\left(y+2\right)^2+\left(z-1\right)^2=9$. Các mệnh đề sau đúng hay sai ?
	\choiceTF
	{\True Thể tích của $\left(S\right)$ là $V=36\pi$ }
	{Phương trình mặt cầu nhận $AB$ làm đường kính có tâm là $I \left(\dfrac{1}{2};1;\dfrac{3}{2}\right)$} 
	{Gọi $K$ là tâm của $\left( S \right)$, khi đó $\mathrm{d}\left(K;\left(ABC\right)\right)=\dfrac{9}{7}$}
	{Tập hợp các điểm $M$ thỏa mãn $MA^2=MB^2+MC^2$ là mặt cầu có bán kính bằng $2$}
	\loigiai{
		\begin{itemchoice}
			\itemch Đúng. Với $R=3$. Thể tích của khối cầu $V = \dfrac{4}{3}\pi R^3=\dfrac{4}{3}\pi 3^3=36\pi$
			\itemch Sai. Phương trình mặt cầu nhận trung điểm $I$ của $AB$ làm tâm. Khi đó $I\left(\dfrac{1}{2};1;0\right)$
			\itemch Sai. Ta có $K\left(3;-2;1\right)$\\
			và phương trình mặt phẳng $\left(ABC\right): \dfrac{x}{1}+\dfrac{y}{2}+\dfrac{z}{3}=1$.\\
			Khi đó, $\mathrm{d}\left(K;\left(ABC\right)\right)=\dfrac{\left\vert \dfrac{3}{1}+\dfrac{-2}{2}+\dfrac{1}{3}-1\right\vert}{\sqrt{\left(\dfrac{1}{1}\right)^2+\left(\dfrac{1}{2}\right)^2+\left(\dfrac{1}{3}\right)^2}}=\dfrac{8}{7}$.
			\itemch Sai. Giả sử $M \left(x;y;z\right)$\\
			Ta có $MA^2=\left(x-1\right)^2+y^2+z^2$; $MB^2=x^2+\left(y-2\right)^2+z^2$; $MC^2=x^2+y^2+\left(z-3\right)^2$.\\
		Khi đó, $MA^2=MB^2+MC^2$ \\
		$ \Leftrightarrow \left(x-1\right)^2+y^2+z^2= x^2+\left(y-2\right)^2+z^2 +x^2+y^2+\left(z-3\right)^2 $ \\
		$\Leftrightarrow -2x+1 = \left( y-2 \right)^2+x^2+\left( y-3\right)^2$\\
		 $\Leftrightarrow \left(x+1\right)^2+\left(y-2\right)^2+\left(z-3\right)^2=2$\\
		Vậy tập hợp các điểm $M$ là mặt cầu có bán kính là $R= \sqrt{2}$.
		\end{itemchoice}
	}
\end{ex}
\begin{ex}%[2H5V3-3]
	Trong không gian với hệ trục tọa độ $Oxyz$ cho mặt phẳng $\left(P\right): 2x+2y-z+16=0$ và mặt cầu $\left(S\right): \left(x-2\right)^2+\left(y-1\right)^2+\left(z-3\right)^2=9$. Các khẳng định nào dưới đây đúng hay sai?
	\choiceTF
	{Khoảng cách từ tâm $I$ của $\left(S\right)$ đến mặt phẳng $\left(P\right)$ bằng $3$} 
	{\True $\left(P\right)$ và $\left(S\right)$ không có điểm chung}
	{Cho mặt phẳng $\left(Q\right):3x-2y+z-10=0$, khi đó $\left(S\right)$ cắt $\left(Q\right)$ theo một giao tuyến là đường tròn có chu vi bằng $\dfrac{3\sqrt{182}}{14} \pi$}
	{\True Cho điểm $M$ di động trên $\left(S\right)$ và điểm $\left(N\right)$ di động trên $\left(P\right)$ sao cho độ dài $MN$ ngắn nhất. Tọa độ điểm $M$ là $M\left(0;-3;4\right)$}
	\loigiai{
		\begin{itemchoice}
			\itemch Sai. Ta có mặt cầu tâm $I\left(2;-1;3\right)$ và có $R=3$. Khoảng cách từ $I$ đến $\left(P\right)$ là\\
			$\mathrm{d}\left(I,\left(P\right)\right)=\dfrac{ \left\vert 2\cdot2+2\cdot\left(-1\right)-3+16 \right\vert}{3}=5$. 
			\itemch Đúng. Vì $\mathrm{d}\left(I,\left(P\right)\right)=5>R=3$ nên $\left(P\right)$ và $\left(S\right)$ không có điểm chung.
			\itemch Sai. Ta có $\mathrm{d}=\mathrm{d}\left(I,\left(Q\right)\right)=\dfrac{3\sqrt{14}}{14}$.\\
			Gọi $r$ là bán kính của đường tròn giao tuyến của $\left(S\right)$ cắt $\left(Q\right)$. \\
			Khi đó $r=\sqrt{R^2-d^2}= \dfrac{3\sqrt{182}}{14}$. Khi đó chu vi cần tìm là $C=2\pi r= \dfrac{3\sqrt{182}}{7} \pi$
			\itemch Đúng. Do $\left(P\right)$ và $\left(S\right)$ không có điểm chung nên $MN_{min}=d-R=2$;\\
			Khi đó $N$ là hình chiếu vuông góc của $I$ lên mặt phẳng $\left(P\right)$ và $M$ là giao điểm của đoạn thẳng $IN$ với mặt cầu $\left(S\right)$\\
			Gọi $\left(d\right)$ là đường thẳng đi qua điểm $I$ và vuông góc với $\left(P\right)$, thì $\left(d\right) \cap \left(P\right) = N$\\
			và $\overrightarrow{u}_{\left(d\right)}=\overrightarrow{u}_{\left(P\right)}=\left(2;2;-1\right)$. Phương trình đường thẳng $\left(d\right): \heva{&x=2+2t\\&y=-1+2t\\&z=3-t}$\\
			$N\in \left(d\right) \Rightarrow N\left(2+2t;-1+2t;3-t\right)$ mà $N \in \left(P\right)$\\
			$\Rightarrow 2\left(2+2t\right)+2\left(-1+2t\right)-\left(3-t\right)+16=0 \Leftrightarrow 9t+15=0 \Leftrightarrow t=\dfrac{-5}{3}$ \\
			Suy ra $N\left(-\dfrac{4}{3};-\dfrac{13}{3};\dfrac{14}{3}\right)$. Ta có $\overrightarrow{IM}=\dfrac{3}{5} \overrightarrow{IN} \Rightarrow M \left(0;-3;4 \right)$.
		\end{itemchoice}
	}
\end{ex}
\begin{ex}%[2H5V3-3]
	Cho mặt cầu $\left(S\right): x^2+y^2+z^2-2x-4y-6z-m=0$.
	\choiceTF
	{Điều kiện của $m$ để mặt cầu $\left(S\right)$ xác định là $m>14$}
	{\True Với $m=2$, đường tròn giao tuyến của $\left(S\right)$ khi cắt mặt phẳng $\left(Oxy\right)$ có chu vi bằng $2\sqrt{7}\pi$}
	{Với $m=2$, Phương trình mặt cầu đối xứng với mặt cầu $\left(S\right)$ qua trục $Oz$ có phương trình $\left(x+1\right)^2+\left(y+2\right)^2+\left(z+3\right)^2=16$}
	{\True Với $m=2$, cho mặt phẳng $\left(P\right): x+y-2x+4=0$. Phương trình đường thẳng $d$ tiếp xúc với $\left(S\right)$ tại $A\left(-1;0;0\right)$ và song song với mặt phẳng $\left(P\right)$ là $\heva{&x=-1+7t\\&y=-7t\\&z=0}$}
	\loigiai{
		\begin{itemchoice}
			\itemch Sai. Phương trình mặt cầu $\left(S\right): x^2+y^2+z^2-2ax-2by-2cz+d=0$ xác định khi $a^2+b^2+c^2-d>0 \Leftrightarrow 14+m>0 \Leftrightarrow m>-14$.
			\itemch Đúng. Với $m=2$, mặt cầu $\left(S\right)$ có tâm $I\left(1;2;3\right)$ và $R=\sqrt{1^2+2^2+3^2+2}=4$. \\
			Ta có $\mathrm{d}\left(I;\left(Oxy\right)\right)=\vert z_I \vert=3$.\\
			Gọi $r$ là bán kính của đường tròn giao tuyến của $\left(S\right)$ và mặt phẳng $Oxy$, ta suy ra $r=\sqrt{R^2-d^2}=\sqrt{7}$. Vậy chu vi cần tìm là $C=2\sqrt{7}\pi$.
			\itemch Sai. Gọi mặt cầu $\left(S'\right)$ là mặt cầu cần tìm. Do $\left(S'\right)$ đối xứng với $\left(S\right)$ qua trục $Oz$ nên tâm $I'$ của $\left(S'\right)$ cũng đối xứng với tâm $I$ của $\left(S\right)$ qua trục $Oz$. Khi đó ta có bán kính $R'=R=4$ và $I'\left(-1;-2;3\right)$. Vậy $\left(S'\right):\left(x+1\right)^2+\left(y+2\right)^2+\left(z-3\right)^2=16 $. 
			\itemch Đúng. Mặt cầu $\left(S\right)$ có tâm $I\left(1;2;3\right)$ nên $\overrightarrow{IA}=\left(-2;-2;-3\right)$. \\
			Đường thẳng $d$ tiếp xúc với mặt cầu $\left(S\right)$ tại $A$ và song song với $\left(P\right)$ nên đường thẳng $d$ có VTCP là $\overrightarrow{u_d}=\left[ \overrightarrow{n}_{\left(P\right)},\overrightarrow{IA}\right] = \left(7;-7;0\right)$.\\
			 Vậy phương trình đường thẳng $d:\heva{&x=-1+7t\\&y=-7t\\&z=0}$ 
		\end{itemchoice}
	}
\end{ex}
\Closesolutionfile{ans}
\indapan{3}{ans/ans-0-B15-DS}

\Opensolutionfile{ans}[ans/ans-0-B15-KQ]
\caukq
\begin{ex}%[2H5H3-1]
	Gọi $I$ là tâm mặt cầu $\left(S \right): x^2+y^2+\left( z-2\right)^2=4$. Độ dài $\left\vert\overrightarrow{OI}\right\vert$ bằng
	\shortans{$2$}
	\loigiai{
		Ta có $I\left(0;0;2\right)$ nên $\overrightarrow{OI}=\left(0;0;2\right)$. Khi đó $\left\vert\overrightarrow{OI}\right\vert = \sqrt{0^2+0^2+2^2}=2$.
	}
\end{ex}

\begin{ex}%[2H5H3-3]
	 Cho mặt cầu $\left(S\right): x^2+y^2+z^2-2x-4y-4z=0$. Giả sử mặt phẳng tiếp xúc với mặt cầu $\left(S\right)$ tại điểm $A\left(3;4;3\right)$ có dạng $ax+by+cz-17=0$. Khi đó $H=a+b+c$ bằng
	\shortans{$5$}
	\loigiai{
		Ta có mặt cầu $\left(S\right)$ có tâm $I\left(1;2;2\right)$. Gọi $\left(\alpha\right)$ là mặt phẳng tiếp xúc với mặt cầu tại điểm $A$, khi đó $\overrightarrow{n}_{\alpha} = \overrightarrow{IA}=\left(2;2;1\right)$\\
		Vậy phương trình mặt phẳng cần tìm là: $$2\left(x-3\right)+2\left(y-4\right)+\left(z-3\right)=0 \Leftrightarrow 2x+2y+z-17=0$$
		Do đó $H=2+2+1=5$
	}
\end{ex}
\begin{ex}%[2H5V3-2]
	Trong không gian với hệ trục tọa độ $Oxyz$, cho $A\left(1;2;3\right)$; $B\left(4;2;3\right)$; $C\left(4;5;3\right)$. Giả sử diện tích mặt cầu nhận đường tròn ngoại tiếp tam giác $ABC$ làm đường tròn lớn có dạng $m\pi$, Khi đó $m$ bằng
	\shortans{$18$}
	\loigiai{
		Ta có $AB = 3$; $BC=3$; $AC=3\sqrt{2}$ nên tam giác $ABC$ vuông cân tại $B$.\\
		Khi đó bán kính đường tròn ngoại tiếp tam giác $ABC$ là $R=\dfrac{3\sqrt{2}}{2}$.\\
		Vậy, diện tích mặt cầu cần tìm là $S=4\pi R^2=18\pi$.
	}
\end{ex}
\begin{ex}%[2H5V3-3]
	Trong không gian với hệ trục tọa độ $Oxyz$, có tất cả bao nhiêu số tự nhiên của tham số $m$ để phương trình $x^2+y^2+z^2+2\left(m-2\right)y-2\left(m+3\right)z+3m^2+7=0$ là phương trình mặt cầu
	\shortans{$4$}
	\loigiai{Theo lí thuyết, Phương trình mặt cầu $\left(S\right): x^2+y^2+z^2-2ax-2by-2cz+d=0$ xác định khi $a^2+b^2+c^2-d>0$. Khi đó, để phương trình $x^2+y^2+z^2+2\left(m-2\right)y-2\left(m+3\right)z+3m^2+7=0$ là phương trình mặt cầu thì $\left(2-m\right)^2+\left(m+3\right)^2-3m^2-7>0$\\
	 $\Leftrightarrow -m^2+2m+6>0 \Leftrightarrow -\sqrt{7}<m<1+\sqrt{7}$, mà $m \in \mathbb{N}$ nên $m \in \left\lbrace 0;1;2;3 \right\rbrace$.
	}
\end{ex}
\begin{ex}%[2H5C3-3]
	Trong không gian $Oxyz$, cho mặt cầu $x^2+y^2+z^2=9$ và điểm $M\left(x_0;y_0;z_0\right)$ thuộc $d: \heva{&x=1+t\\&1+2t\\&2-3t}$. Ba điểm $A$,$B$,$C$ phân biệt cùng thuộc mặt cầu sao cho $MA$,$MB$,$MC$ là các tiếp tuyến của mặt cầu. Biết rằng mặt phẳng $\left(ABC\right)$ đi qua $D\left(1;1;2\right)$. Tổng $T=x_{0}^2+y_{0}^2+z_{0}^2$ bằng
	\shortans{$26$}
	\loigiai{
	Ta có mặt cầu $\left(S\right)$ có tâm $O\left(0;0;0\right)$ và $R=3$. Gọi $M\left(1+t_0;1+2t_0;2-3t_0\right) \in d$.\\
	Giả sử $T\left(x;y;z\right) \in \left(S\right)$ là một tiếp điểm của tiếp tuyến $MT$ với mặt cầu $\left(S\right)$. Khi đó $OT^2+MT^2=OM^2$\\
	$\Leftrightarrow 9+\left[x-\left(1+t_0\right)\right]^2+\left[y-\left(1+2t_0\right)\right]^2+\left[z-\left(2-3t_0\right)\right]^2=\left(1+t_0\right)^2+\left(1+2t_0\right)^2+\left(2-3t_0\right)^2$\\
	$\Leftrightarrow \left(1+t_0\right)x+\left(1+2t_0\right)y+\left(2-3t_0\right)z-9=0$.\\
	Suy ra phương trình mặt phẳng $\left(ABC\right)$ có dạng $\left(1+t_0\right)x+\left(1+2t_0\right)+\left(2-3t_0\right)z-9=0$. Hơn nữa $D \in \left(ABC\right)$ nên thay $x=1$, $y=1$, $z=2$ vào $\left(ABC\right)$ ta được $t_0=-1$\\
	$\Rightarrow M\left(0;-1;5\right)$. Vậy $T=0^2+\left(-1\right)^2+5^2=26$
	}
\end{ex}
\begin{ex}%[2H5C3-3]
	Trong không gian $Oxyz$, cho các điểm $M\left(2;1;4\right)$, $N\left(5;0;0\right)$, $P\left(1;-3;1\right)$. Gọi $I\left(a;b;c\right)$ là tâm của mặt cầu tiếp xúc với mặt phẳng $\left(Oyz\right)$ đồng thời đi qua các điểm $M,N,P$. Biết rằng $a+b+c<5$. Khi đó $c$ bằng 
	\shortans{$2$}
	\loigiai{
		Phương trình mặt cầu $\left(S\right)$ tâm $I\left(a;b;c\right)$ có dạng $x^2+y^2+z^2-2ax-2by-2cz+d=0$, với điều kiện $a^2+b^2+c^2-d>0$\\
		Do $\left(S\right)$ đi qua các điểm $M,N,P$ và tiếp xúc với mặt phẳng $\left(Oyz\right)$ nên ta có\\
		$\Leftrightarrow \heva{&-4a-2b-8c+d=-21\\&-10a+d=-25\\&-2a+6b-2c+d=-11\\&R=\vert a \vert} \Leftrightarrow \heva{&-4a-2b-8c+10a-25=-21\\&d=10a-25\\&-2a+6b-2c+10a-25=-11\\&a^2+b^2+c^2-d=a^2}$\\
		$\Leftrightarrow \heva{&6a-2b-8c=4\\&d=10a-25\\&8a+6b-2c=14\\&b^2+c^2-d=0} \Leftrightarrow \heva{&6a-2b-8c=4\\&d=10a-25\\&32a+24b-8c=52\\&b^2+c^2-d=0} \Leftrightarrow \heva{&6a-2b-8c=4\\&d=10a-25\\&26a+26b=52\\&b^2+c^2-d=0}$\\
		$\Leftrightarrow \heva{&c=a-1\\&d=10a-25\\&b=-a+2\\&b^2+c^2-d=0}$\\
		$\Leftrightarrow \left(-a+2\right)^2+\left(a-1\right)^2-10a+25=0 \Leftrightarrow 2a^2-16a+30=0$\\
		$\Leftrightarrow \hoac{&a=3\\&a=5}$. Do $a+b+c<5$ nên $a=3,b=-1, c=2$
	}
\end{ex}
\Closesolutionfile{ans}
\indapan{6}{ans/ans-0-B15-KQ}
